\documentclass[]{article}
\usepackage{graphicx}
\usepackage[spanish]{babel}
\usepackage[a4paper, top=2.5cm, bottom=2.5cm, left=3cm, right=3cm]{geometry}
\usepackage[hidelinks]{hyperref}
\usepackage[T1]{fontenc}
\usepackage{listings}
\usepackage{xcolor}
\usepackage{float}

\definecolor{miverde}{rgb}{0,0.6,0}
\lstdefinelanguage{flag}{
  keywords={dbafcc, cafebb},
  keywordstyle=\bfseries,
  ndkeywords={},
  ndkeywordstyle=\color{miotrootroverde}\bfseries,
  identifierstyle={},
  sensitive=false,
  comment=[l]{//},
  morecomment=[s]{/*}{*/},
  commentstyle=\color{red}\ttfamily,
  stringstyle=\color{miverde}\ttfamily,
  morestring=[b]',
  morestring=[b]"
}
% style for listings (código)
\lstdefinestyle{python}{
    language=Python,
    backgroundcolor=\color{gray!2},     % Color de fondo
    basicstyle=\ttfamily,               % Tipo y tamaño de fuente
    keywordstyle=\color{blue}\bfseries, % Color para palabras clave
    stringstyle=\color{miverde},        % Color para cadenas
    commentstyle=\color{red},           % Color para comentarios
    showspaces=false,                   % No mostrar espacios
    showstringspaces=false,             % No mostrar espacios en las cadenas
    frame=single,                       % Poner un marco alrededor del código
    breaklines=true,                    % Romper las líneas largas
    captionpos=b,                       % Posición del caption
    tabsize=4,                          % Tamaño de las tabulaciones
    escapeinside={\%*}{*)},             % Para incluir código LaTeX en los listings
    morekeywords={self}                 % Palabras clave adicionales
}

\lstdefinestyle{flag}{
    language=bash,
    backgroundcolor=\color{gray!2},     % Color de fondo
    basicstyle=\ttfamily,               % Tipo y tamaño de fuente
    keywordstyle=\bfseries, % Color para palabras clave
    stringstyle=\color{miverde},        % Color para cadenas
    commentstyle=\color{red},           % Color para comentarios
    showspaces=false,                   % No mostrar espacios
    showstringspaces=false,             % No mostrar espacios en las cadenas
    frame=single,                       % Poner un marco alrededor del código
    breaklines=true,                    % Romper las líneas largas
    captionpos=b,                       % Posición del caption
    tabsize=4,                          % Tamaño de las tabulaciones
    escapeinside={\%*}{*)},             % Para incluir código LaTeX en los listings
    morekeywords={self}                 % Palabras clave adicionales
}


\lstdefinestyle{bash}{
    language=bash,
    backgroundcolor=\color{gray!2},     % Color de fondo
    basicstyle=\ttfamily,               % Tipo y tamaño de fuente
    keywordstyle=\bfseries, % Color para palabras clave
    stringstyle=\color{miverde},        % Color para cadenas
    commentstyle=\color{red},           % Color para comentarios
    showspaces=false,                   % No mostrar espacios
    showstringspaces=false,             % No mostrar espacios en las cadenas
    frame=single,                       % Poner un marco alrededor del código
    breaklines=true,                    % Romper las líneas largas
    captionpos=b,                       % Posición del caption
    tabsize=4,                          % Tamaño de las tabulaciones
    escapeinside={\%*}{*)},             % Para incluir código LaTeX en los listings
    morekeywords={self}                 % Palabras clave adicionales
}
\lstset{basicstyle=\ttfamily}
\hypersetup{
    colorlinks=true, % Activa el color en los enlaces
    linkcolor=blue,  % Color de los enlaces internos (como secciones)
    urlcolor=blue,   % Color de los enlaces a URLs
}


%title
\title{Práctica 1} 

\author{Adrián Ferández Galán, César López Mantecón y Manuel Gómez-Plana Rodríguez}

\begin{document}

\begin{titlepage}
    \centering
   \includegraphics[width=0.9\textwidth]{uc3m.jpg} 
    {\Huge Universidad Carlos III\\
    
     \Large Ingeniería de la Ciberseguridad\\
     \vspace{0.5cm}
     Curso 2024-25}
    \vspace{2cm}

    {\Huge \textbf{Práctica 2} \par}
    \vspace{0.5cm}
    {\Large Intercepción de tráfico HTTP\par}
    \vspace{8cm}

   \textbf{Ingeniería Informática, Cuarto curso}\\
    \vspace{0.2cm} 
    Adrián Fernández Galán       (NIA: 100472182, e-mail: 100472182@alumnos.uc3m.es)\\
    César López Mantecón         (NIA: 100472092, e-mail: 100472092@alumnos.uc3m.es)\\
    Manuel Gómez-Plana Rodríguez (NIA: 100472310, e-mail: 100472310@alumnos.uc3m.es)
    \vspace{0.5cm}

   
    \textbf{Prof.} Antonio Nappa\\
    
    \textbf{Grupo: } 81   
    
\end{titlepage}
\newpage

%\renewcommand{\contentsname}{\centering Índice}
%\tableofcontents

\newpage

\lstset{style=bash}
\section{Introducción}
\label{sec:introduccion}
Este documento recoge el desarrollo de la segunda práctica de la asignatura
\textit{Ingeniería de la Ciberseguridad}. El objetivo de esta práctica es
desencriptar una serie de \textit{flags} contenidas en paquetes \texttt{http}
que atraviesan un \textit{proxy}.

Exiten dos clases de \textit{flags} contenidas en los paquetes:
\begin{itemize}
    \item \textbf{Tipo 1}: se trata de una cadena hexadecimal con un argumento
    llamado \textit{rotation}.
    \item \textbf{Tipo 2}: se trata de una cadena de caracteres con 3
    argumentos (\textit{Rotation, s} y \textit{op})
\end{itemize}

\section{Funcionamiento de \texttt{mitmproxy}}
\label{sec:proxy}

\texttt{Mitmproxy} es una herramienta para la consola de comandos que permite
el análisis interactivo y la modificaicón de tráfico \texttt{http}. Por
defecto, el proxy esucha sobre la dirección \textit{localhost} y puerto
\textit{8080}.

Adicionalmente, existe un proceso en constante ejecución llamado
\textit{down.sh} que hace una petición http sobre esta misma dirección.

\begin{lstlisting}[language=bash, caption=Resultado de ejecutar ps -e]
   PID TTY          TIME CMD
      1 pts/0    00:00:00 bash
    587 ?        00:00:00 sshd
  26023 ?        00:00:37 down.sh
 900981 ?        00:00:00 sshd
 900998 ?        00:00:00 sshd
 900999 pts/1    00:00:00 bash
 901270 ?        00:00:00 sleep
 901273 pts/1    00:00:00 ps
\end{lstlisting}

Se ha utilizado el comando \texttt{find} para localizar el archivo. Aprovechando el conocimiento previo que tenemos sobre la máquina, hemos empleado el usuario \textit{admin} para poder lanzar el comando sobre la raíz, ya que este usuario pertenece al grupo de superusuarios: 

\begin{lstlisting}[language=bash, caption=Resultado de ejecutar find]
sudo find / -type f -name down.sh
[sudo] password for admin:
./down.sh
\end{lstlisting}

Finalmente, analizando el contenido del archivo podemos ver el comando
utilizado para realizar la petición \texttt{http}.
\begin{lstlisting}[language=bash, caption=petición http a través del comando curl]
    /home/rover/curl-7.88.1/src/curl --insecure --proxy-insecure --proxy https://172.17.0.19:8080 https://172.17.0.4:4443/
\end{lstlisting}

Se observa que el comando redirige su petición a la ip y puerto donde escucha
\texttt{mitmproxy}. Adicionalmente, utiliza las flags \texttt{--insecure} y
\texttt{--proxy-insecure}, por lo que no verificará los certificados del proxy
y se permitirá la conexión. Con esto, la conexión se completa sin levantar
excepciones y podremos ser capaces de analizar las peticiones y respuestas a
través de la interfaz de \texttt{mitmproxy}.

\section{Descubrimiento de una flag tipo 1}
\label{sec:tipo1}

Las \textit{flags} están codificadas en hexadecimal y cuentan con un parámetro
llamado \textit{Rotation}. Este parámetro parece indicar que se ha cifrado a
través de una rotación en la cadena de texto. Al analizar varios paquetes se
aprecia que las flags de igual rotación cuentan con caracteres entre los
índices 10 y 15, lo que sugiere la utilización de una misma semilla para todas
las flags.

\lstset{style=flag}
\begin{lstlisting}[language=flag, caption=Varias flags para rotation 1 donde la semilla es dbafcc]
3193601502dbafcc86bdb57381081aec
3846408208dbafccc0399a91c4fbf1a3
4745131622dbafcc09340e7b0af6885f
\end{lstlisting}

Primeramente se trató de desencriptar las \textit{flags} rotando en ambos sentidos los valores hexadecimales de cada caracter. No obstante, esto no dio resultado. Adicionalmente, se probaron otras codificaciones como \texttt{ascii}.

Posteriormente se observó que las flags con una diferencia de 6 rotaciones mantenían la misma semilla. Dado que la semilla está formada enteramente por letras y que la codificación hexadecimal cuenta con 6 caracteres alfabéticos, esto sugiere que el cifrado se ha aplicado por separado a letras y números.

\begin{lstlisting}[language=flag, escapechar=@, caption=Varias flag con la misma semilla y 6 rotaciones]
FLAG                              Rot
3193601502dbafcc86bdb57381081aec    1
3846408208dbafccc0399a91c4fbf1a3    1
4745131622dbafcc09340e7b0af6885f    1
9352846549dbafcc965695c25417c6d9    7
7173993205dbafccf2c4f4db398e6687    7
1047974726dbafcc05aacde26d253963    7
7512913585dbafcc9aeb0520366299c2    7
5483152252dbafcc26492c301ca7365e    7
2473284585dbafcc99892775f126bd30   13
0872242899dbafccfc622e537f7ba143   13
9735635169dbafcc24d5a92548bc1e09   13
7646156765dbafccd0beba631d833c5c   13
\end{lstlisting}

Finalmente, se ha hecho la prueba obteniendo una \textit{flag} válida a través
de un programa de python.

\lstset{style=python}
\begin{lstlisting}[language=Python, caption=Decodificación de una flag de rotación 1]
flag = "9443878358dbafcc223a8e670ba03537"
num_rotations = 1
letters_to_number = {'a': 0, 'b': 1, 'c': 2, 'd': 3, 'e': 4, 'f': 5}
number_to_letters = {0: 'a', 1: 'b', 2: 'c', 3: 'd', 4: 'e', 5: 'f'}
new_flag = ""

for i in range(len(flag)):
    c = flag[i]
    if c.isnumeric():
        new_c = (int(c) - num_rotations) % 10
    elif c in letters_to_number.keys():
        new_c = number_to_letters[(letters_to_number[c] - num_rotations) % 6]
    else:
        new_c = c
    new_flag += str(new_c)

print(new_flag)

\end{lstlisting}

El resultado obtenido ha sido el siguiente:
\texttt{8332767247cafebb112f7d569af92426}. Este resultado si ha sido aceptado
como correcto.

A modo de curiosidad, dado que el cifrado de rotación es modular, la semilla se
mostrará en claro para una rotación igual al periodo. Es por esto que se puede
apreciar la semilla en las flags con rotación 6:

\begin{lstlisting}[language=flag, caption=Flags para rotacion 6]
3786566793cafebb39330ab797fbecfa
6350772345cafebb35f756efb7386a9d
0706541085cafebbeaf7963399a4d327
2847807583cafebbd1a054ded7579a71
1613995163cafebb528dd12a5f194769
\end{lstlisting}

\section{Descubrimiento de una flag tipo 2}
\label{sec:tipo2}

\section{Conclusión}
\label{sec:conclusion}

\end{document}
